\section{四类序列模型回答不同问题}
\label{sec:13-Machine-Learning/07-Graphical-and-Sequential-Models/05}

设备故障预警的选型会已经开了两个小时。手上的数据是每分钟一条的多路传感器序列，外加运维人员事后标出的少量故障区间。一位同事主张用隐 Markov 模型，说状态可解释；另一位要上 Transformer，理由是长程依赖；还有人提议 LSTM 接一层 CRF，因为标注是区间。争论一直卡着，因为大家在比``哪个更强''。可这四样东西建模的对象不同，连随机变量是谁都不一样，强弱无从比起。

\subsection{同一条链，条件线画在不同位置}

先把四类模型的图画出来，看谁被观察、谁是随机的。

\begin{center}
\begin{tikzpicture}[>=stealth,line width=0.7pt]
% ---------- HMM ----------
\begin{scope}[shift={(0,2.9)}]
  \foreach \x in {0,1.3,2.6}{
    \fill[black!15] (\x,0) circle (0.26);
    \draw (\x,0) circle (0.26);
    \draw (\x,1.15) circle (0.26);
    \draw[->] (\x,0.89) -- (\x,0.26);
  }
  \draw[->] (0.26,1.15) -- (1.04,1.15);
  \draw[->] (1.56,1.15) -- (2.34,1.15);
  \node[font=\small] at (0,1.15) {\(Z_1\)};
  \node[font=\small] at (1.3,1.15) {\(Z_2\)};
  \node[font=\small] at (2.6,1.15) {\(Z_3\)};
  \node[font=\small] at (0,0) {\(X_1\)};
  \node[font=\small] at (1.3,0) {\(X_2\)};
  \node[font=\small] at (2.6,0) {\(X_3\)};
  \node[anchor=west,font=\small] at (-0.3,1.95) {HMM：离散隐状态，联合分布};
\end{scope}
% ---------- 线性 Gauss 状态空间 ----------
\begin{scope}[shift={(6.2,2.9)}]
  \foreach \x in {0,1.3,2.6}{
    \fill[black!15] (\x,0) circle (0.26);
    \draw (\x,0) circle (0.26);
    \draw (\x,1.15) circle (0.26);
    \draw[->] (\x,0.89) -- (\x,0.26);
  }
  \draw[->] (0.26,1.15) -- (1.04,1.15);
  \draw[->] (1.56,1.15) -- (2.34,1.15);
  \node[font=\small] at (0,1.15) {\(Z_1\)};
  \node[font=\small] at (1.3,1.15) {\(Z_2\)};
  \node[font=\small] at (2.6,1.15) {\(Z_3\)};
  \node[font=\small] at (0,0) {\(X_1\)};
  \node[font=\small] at (1.3,0) {\(X_2\)};
  \node[font=\small] at (2.6,0) {\(X_3\)};
  \node[anchor=west,font=\small] at (-0.3,1.95) {线性 Gauss：同图，状态连续};
\end{scope}
% ---------- 线性链 CRF ----------
\begin{scope}[shift={(0,0)}]
  \foreach \x in {0,1.3,2.6}{
    \fill[black!15] (\x,0) circle (0.26);
    \draw (\x,0) circle (0.26);
    \draw (\x,1.15) circle (0.26);
  }
  \draw (0.26,1.15) -- (1.04,1.15);
  \draw (1.56,1.15) -- (2.34,1.15);
  \foreach \a in {0,1.3,2.6}{
    \foreach \b in {0,1.3,2.6}{
      \draw[black!30] (\a,0.26) -- (\b,0.89);
    }
  }
  \node[font=\small] at (0,1.15) {\(Y_1\)};
  \node[font=\small] at (1.3,1.15) {\(Y_2\)};
  \node[font=\small] at (2.6,1.15) {\(Y_3\)};
  \node[font=\small] at (0,0) {\(X_1\)};
  \node[font=\small] at (1.3,0) {\(X_2\)};
  \node[font=\small] at (2.6,0) {\(X_3\)};
  \node[anchor=west,font=\small] at (-0.3,1.95) {CRF：输入全部在条件线右边};
\end{scope}
% ---------- self-attention ----------
\begin{scope}[shift={(6.2,0)}]
  \foreach \x in {0,1.3,2.6}{
    \fill[black!15] (\x,0) circle (0.26);
    \draw (\x,0) circle (0.26);
    \draw (\x-0.26,0.89) rectangle (\x+0.26,1.41);
  }
  \foreach \a in {0,1.3,2.6}{
    \foreach \b in {0,1.3,2.6}{
      \draw[->,black!35] (\a,0.26) -- (\b,0.87);
    }
  }
  \node[font=\small] at (0,1.15) {\(h_1\)};
  \node[font=\small] at (1.3,1.15) {\(h_2\)};
  \node[font=\small] at (2.6,1.15) {\(h_3\)};
  \node[font=\small] at (0,0) {\(X_1\)};
  \node[font=\small] at (1.3,0) {\(X_2\)};
  \node[font=\small] at (2.6,0) {\(X_3\)};
  \node[anchor=west,font=\small] at (-0.3,1.95) {注意力：确定性映射，无后验};
\end{scope}
\end{tikzpicture}
\end{center}

\emph{灰底表示被观察；方框表示确定性中间量，不是随机变量。}

左上与右上两张图完全一样，差别只在状态取值是有限类别还是连续向量——正是这一点决定了前向递推里的求和变积分、概率表变协方差。左下的 CRF 把输入整段搬到条件线右边，于是允许任意重叠的特征接到标签上，换来这份自由，它不再拥有 \(p(x)\)。右下的注意力层里没有一个圆圈画在上排，那些方框是输入的确定性函数，谈不上后验。

\subsection{建模对象决定了能问哪些问题}

把上图翻译成一张对照表。

\begin{center}
\begin{tabular}{p{0.17\textwidth}p{0.20\textwidth}p{0.26\textwidth}p{0.24\textwidth}}
\toprule
模型 & 主要分布或映射 & 结构假设 & 典型问题 \\
\midrule
HMM &
\(p(x_{1:T},z_{1:T})\) &
离散一阶 Markov 状态；给定状态后观测独立 &
状态解释、生成、滤波与平滑 \\
线性 Gauss 状态空间 &
\(p(x_{1:T},z_{1:T})\) &
线性动力学与 Gauss 噪声 &
跟踪、外推、带不确定性的状态估计 \\
线性链 CRF &
\(p(y_{1:T}\given x_{1:T})\) &
给定整段输入后，标签图为链 &
结构化标注与全局一致解码 \\
自注意力 &
\(h_{1:T}=F_\theta(x_{1:T})\) &
内容相关的两两读取；概率语义由损失另行规定 &
上下文表示，作为其他模型的组件 \\
\bottomrule
\end{tabular}
\end{center}

第二列最要紧。写了联合分布的两行可以谈缺失观测怎么填、未来怎么生成、状态的后验有多宽；只写条件分布的那行避开了为高维输入建生成模型的负担，但也就没法比较两段输入谁更可能出现；只有映射的那行要接上 softmax 似然、自回归分解或潜变量结构之后，才谈得上概率。

回到故障预警。若运维想问``设备现在处于哪个退化阶段、把握有多大''，需要状态的后验，前两行才回答得了。若只想把已经录好的历史序列切成正常段与故障段，第三行更省事。若手上标注很少而原始序列很多，先用第四行学表示、再接一个概率头，往往比强行给传感器写生成模型更实际。争论谁更强之前，先说清楚要问的是哪一类问题。

\subsection{同叫状态，语义可能完全不同}

HMM 的 \(Z_t\) 与 Kalman filter 的 \(X_t\) 是随机潜变量，看完数据仍留有后验不确定性，你可以问它取某个值的概率是多少。循环网络或注意力层里的隐向量是输入与参数的确定性函数，把它叫做``隐状态''不等于定义了 \(p(z_t\given x_{1:t})\)，向量的范数也不能当置信度读。想要后验，得另外规定潜变量与似然，再用变分或采样去近似。

反过来说，有概率语义也不保证预测更好。真实动力学若强非线性、依赖跨度大或状态维数设错，一个漂亮的一阶 Markov 模型会给出系统性偏差，而且它照样输出一个窄窄的协方差——上一单元讲的那件事在这里重现：模型只在自己的假设内部诚实。概率语义让不确定性这个问题可以被提出来，不负责让答案正确。

\subsection{处理长序列不是一种单一能力}

四类模型的算法看似各不相同，实际都在利用可复用的结构。链上的和积递推合并相同前缀，把路径求和压到与长度成正比；Gauss 族在线性条件化下闭合，只需传播均值与协方差；CRF 用前向--后向算条件配分函数、用 Viterbi 取最高分路径；注意力则直接做内容相关的加权和，标准实现的两两交互要 \(O(T^2)\) 的时间与存储。

于是``能处理长序列''分成了两件事。链式动态规划对长度线性，但它只能表达局部因子，远距离依赖必须经由状态一路传递，中途容易被噪声冲淡。注意力一步就能连上任意两个位置，要付的是二次的交互开销，而且它不自带滤波后验——想要在线的状态估计，还得在外面再套一层概率结构，或者退到\hyperref[sec:13-Machine-Learning/08-Sampling-and-Inference/05]{``粒子滤波用带权样本追踪非线性状态''}那类近似方法。注意力机制本身的代价与变体见\hyperref[ch:127]{``注意力与序列模型：读谁交给内容决定，位置、深度与算力另外付账''}。

\subsection{组合模型时要守住接口}

实际系统很少只用一样。上下文网络先给每个位置打分、CRF 层再保证转移合法；神经网络参数化非线性状态转移、粒子滤波近似后验；注意力产生策略的输入表示、策略分布决定动作。组合不消解原来的问题，只是把它们分到三个接口上——表示能力、概率目标、推断算法。

排查失败时沿这三层走一遍通常比换模型有效。训练损失下不去，多半是表示不足或优化出了问题；似然正常而滤波结果漂移，先怀疑状态转移设错或近似退化；逐点得分合理而输出序列非法，说明输出结构压根没被建模。知道一个方法属于哪一层，比记住它``适合序列''更能指导选择。

本单元从头到尾只做了一件事：把条件独立当作可以书写、可以检查的建模资源。它换来的是参数量和计算量的下降，付出的是一组可能出错而不会报错的假设。图画得越干净，越需要有人回头去问那些缺失的边凭什么缺失。
